%---------------------------------------------------------------------
\subsection{Présentation synthétique des responsabilités}
%---------------------------------------------------------------------

Si j'ai exercé une série de responsabilités collectives par le passé (participation à des conseils locaux et nationaux, direction d'équipe, direction d’UMR, etc), je n’en ai exercé formellement aucune depuis mon arrivée à \SU. Je participe néanmoins à différentes instances d’intérêt général.

% %---------------------------------------------------------------------
% \subsection{Responsabilités administratives}
% %---------------------------------------------------------------------
% %     • Présidence ou vice-présidence d'établissement de l'enseignement supérieur
% %     • Direction de composante, d'école doctorale, de services communs
% %     • Direction de structures de recherche ou d’appui (UMR, EA, SFR, ERT, plateformes ...)
% %     • Missions et gestion de projets de l'établissement
% %     • Autres
% J’ai été directeur de l’\UMRlong pendant neuf ans, de 2004 à 2013, période durant laquelle l’unité a reçu trois évaluations de l’AERES, devenue ensuite HCERES. L’unité comptait alors un peu plus de 20 permanents et de 30 membres au total répartis en deux ou trois équipes, selon les périodes, et travaillant sur des sujets variés en informatique, statistique et probabilités appliqués, avec des applications en risque sanitaire, environnement ou génomique. 

% %---------------------------------------------------------------------
% \subsection{Responsabilités et mandats locaux ou régionaux}
% %---------------------------------------------------------------------
% %     • Participation aux conseils centraux (rôle, missions…)
% %     • Participation aux conseils de composantes, de laboratoires...
% %     • Autres
% Pendant mon long séjour à \APT, j’ai été successivement élu dans toutes ses instances : conseil des enseignants, conseil scientifique et conseil d’administration. J’ai également siégé dans les conseils scientifiques de plusieurs départements de l’INRAE.
% 
% J'ai également été membre du comité de pilotage de la fondation mathématique Jacques Hadamard (FMJH) et de l'école doctorale de mathématiques Hadamard (EDMH) lors de sa création. Je suis ensuite devenu membre du conseil scientifique de la FMJH jusqu’en 2022. 

%---------------------------------------------------------------------
\subsection{Responsabilités et mandats (internationaux, nationaux)}
%---------------------------------------------------------------------
% %     • Participations à des instances nationales - CNU, CNRS, conseils des établissements publics, jurys de concours – ou internationales
% %     • Responsabilités exercées dans les agences nationales (HCERES, ANR, ...)
% %     • Autres
% Dans le cadre universitaire, j’ai été nommé au CNU en section 26 où j’ai siégé pendant deux ans. 
% 
% J’ai fait partie du comité AERES pour l’Institut Mathématique de Toulouse en 2014, j’ai dirigé le comité HCERES pour l’Unité de Sensométrie et Chimiométrie (ONIRIS, Nantes) en 2016. J’ai participé à l’évaluation HCERES de l’unité BioSTDM (univ. Paris Descartes) en 2018. 
% J’ai également été rapporteur dans les évaluations de plusieurs unités INRIA (MISTIS en 2007, SISTM en 2014, CELESTE en 2022).

J’ai été rapporteur de l'évaluation de l'unité INRIA (univ. Paris-Saclay) CELESTE.

%---------------------------------------------------------------------
\subsection{Autres informations}
%---------------------------------------------------------------------
Je suis membre du comité de pilotage de la fondation pour la Science Mathématique de Paris centre (FSMP) depuis 2024. Cette instance coordonne les actions de la fondation au fil de l’eau. Je participe notamment au jury d’attribution des bourses de post-doc de FSMP.

Je suis également membre du comité exécutif de l'initiative InLife de \SU (Sciences aux Interfaces du Vivant) qui soutient les actions de recherche à l’interface entre les sciences du vivant et toutes les autres disciplines (physique, chimie, mathématiques, mais aussi sciences humaines et sociales). À ce titre, je participe au jury d’attribution des contrats doctoraux du programme doctoral Interfaces pour le vivant (IPV) de \SU.
